\documentclass[11pt,letterpaper]{article}

\usepackage[margin=1in]{geometry}
\usepackage[T1]{fontenc}
\usepackage{lmodern}
\usepackage{microtype}
\usepackage{parskip}
\usepackage{enumitem}
\usepackage{titlesec}
\usepackage[hidelinks,colorlinks=true,linkcolor=black,urlcolor=black,citecolor=black]{hyperref}
\usepackage{xcolor}
\usepackage{fancyhdr}

\definecolor{theseusbrown}{RGB}{148,110,78}

\titleformat{\section}[block]{\normalsize\scshape\color{theseusbrown}}{}{0em}{}
\titlespacing*{\section}{0pt}{1.8em}{0.4em}

\pagestyle{fancy}
\fancyhf{}
\fancyhead[L]{\small\textsc{\color{theseusbrown}Theseus}}
\fancyhead[R]{\small\textsc{\color{theseusbrown}Questions}}
\fancyfoot[C]{\thepage}
\renewcommand{\headrulewidth}{0pt}

\setlength{\parindent}{0pt}

\begin{document}

\thispagestyle{empty}

\vspace*{2cm}

\begin{center}
{\Large\textsc{\textbf{Eight Questions}}}\\[8pt]
{\normalsize On the Purpose of Schooling}\\[16pt]
{\small\textit{Theseus Intellectual Capital Partners \quad$\cdot$\quad April 2026}}
\end{center}

\vspace{4pt}
\noindent\rule{\textwidth}{0.4pt}
\vspace{0.8em}

\section{On the Conflation of Learning and School}

\textbf{1.}
Every human society ever observed educates its young. Almost none of them, until very recently, did so by confining children to purpose-built rooms for twelve consecutive years. We now treat these two facts---that humans learn, and that humans attend school---as if they describe the same phenomenon. What if the relationship between learning and schooling is more like the relationship between nutrition and the restaurant industry: real, but far more contingent and far less necessary than the people who run restaurants would like you to believe?

\section{On the Sorting Function}

\textbf{2.}
Schools test, rank, grade, and credential. Employers rely on these signals not because they indicate competence---decades of research show the correlation is weak---but because they indicate compliance, endurance, and the willingness to defer judgment to an institution for years on end. If the primary economic function of schooling is to sort people by temperament rather than to develop their capacities, then the students who rebel are not failing the system; the system is confessing what it actually selects for. Who benefits from making this confession inaudible?

\section{On the Hidden Curriculum}

\textbf{3.}
The explicit curriculum teaches mathematics, literature, history. The hidden curriculum teaches something else: how to sit still on command, how to care about a grade more than a question, how to produce work for an audience of one who already knows the answer, how to suppress the urge to speak until called upon. These lessons are never tested, never graded, and never debated at school board meetings---and yet they may be the only lessons that every student, without exception, actually learns. If the hidden curriculum is the real curriculum, what is the explicit one for?

\section{On Curiosity and Its Extinction}

\textbf{4.}
A five-year-old asks roughly forty thousand questions a year. By the time that same child is twelve, the rate has collapsed. The standard explanation is developmental: children mature out of indiscriminate questioning and into focused study. But the collapse is suspiciously well-timed with the first years of formal schooling, and it is suspiciously well-suited to the needs of an institution that must move thirty people through the same material at the same pace. Is the decline of curiosity a natural stage of cognitive development---or an injury we have learned to describe as growth?

\section{On the Identical Classroom}

\textbf{5.}
A photograph of a classroom in 1890 and a photograph of a classroom in 2025 are, structurally, almost indistinguishable: rows of seats, a single authority figure at the front, a board, a bell. In the same period, medicine abandoned humoral theory, physics discovered and harnessed quantum mechanics, and communication went from telegraph to satellite. The classroom's immunity to the forces that transformed every other institution is either evidence that it was already perfect---or evidence that its persistence is explained by something other than effectiveness. Which is more likely, and what would that something be?

\section{On Socialisation as Justification}

\textbf{6.}
When every other argument for schooling is challenged---academic outcomes are uneven, economic returns are overstated, credentialing is circular---the final redoubt is socialisation. Children must go to school to learn how to be around other people. But the model of socialisation that school provides is historically bizarre: age-segregated, authority-mediated, and almost entirely disconnected from the adult world the child will eventually enter. If a sociologist from another civilisation observed the arrangement, would they describe it as preparation for social life---or as a systematic way of preventing children from participating in it?

\section{On What Cannot Be Measured}

\textbf{7.}
Standardised testing emerged from the same early-twentieth-century impulse that gave us scientific management, eugenics, and the assembly line: the conviction that all meaningful variation can be quantified and ranked. The things tests measure---recall, pattern recognition, processing speed under time pressure---are real cognitive abilities. The things they cannot measure---moral seriousness, the capacity to sit with uncertainty, the willingness to change one's mind at real cost---may be more important to every outcome that actually matters. If the instrument determines what counts as achievement, and the instrument is structurally blind to the qualities most worth cultivating, then optimising for test performance is not a path toward excellence. What is it a path toward?

\section{On the Unasked Question}

\textbf{8.}
We debate how to improve schools---better teachers, better funding, better curricula, better technology---with an intensity that makes it almost impossible to ask whether the institution itself is the right unit of analysis. The question \textit{what is school for?} is treated as either trivially obvious or dangerously nihilistic, which is precisely how a society treats questions it cannot afford to answer. What would it mean to take the question seriously---not as an attack on teachers or children, but as a genuine inquiry into whether the thing we have built is the thing we need?

\vspace{2em}
\begin{center}
\rule{3cm}{0.4pt}\\[1em]
{\small\textit{Theseus Intellectual Capital Partners \quad$\cdot$\quad Podcast Discussion Material}}
\end{center}

\end{document}
